\documentclass[11pt, letterpaper]{article}
\usepackage[utf8]{inputenc}
\usepackage[T1]{fontenc}
\usepackage{geometry}
\geometry{top=0.75in, bottom=0.75in, left=1in, right=1in}
\usepackage{enumitem}
\usepackage{amssymb}
\usepackage{hyperref}
\usepackage{titlesec}
\usepackage{color}
\usepackage{fancyhdr}

% Official Columbia Primary Color (Dark Blue)
\definecolor{columbiablue_official}{RGB}{0, 33, 71}

\pagestyle{fancy}
\fancyhf{}
\renewcommand{\headrulewidth}{0.5pt}
\renewcommand{\headrule}{\hbox to\headwidth{\color{columbiablue_official}\leaders\hrule height \headrulewidth\hfill}}

% Header
\lhead{\small \textbf{\color{columbiablue_official} COLUMBIA UNIVERSITY} \\ \href{https://www.arch.columbia.edu/}{Graduate School of Architecture, Planning and Preservation}}
\rhead{\small \href{https://www.linkedin.com/in/dhardestylewis}{Daniel Hardesty Lewis} \\ Ph.D. Candidate, Urban Planning} % Kept Ph.D as per your overview doc, adjust if needed

% Section styling - Clean and branded
\titleformat{\section}{\large\bfseries\color{columbiablue_official}}{}{0em}{}[\color{columbiablue_official}\titlerule]
\titlespacing*{\section}{0pt}{8pt}{3pt}

\begin{document}
\pagestyle{fancy}

\vspace*{0.2cm}

\begin{center}
    {\Large \textbf{\color{columbiablue_official} Thesis Outline}} \\[0.4em]
    {\large \textit{\href{https://github.com/dhardestylewis/thesis}{Predicting Zoning Opposition}}}
\end{center}

\vspace{0.2cm}

\section*{1. Introduction (3-5 pages)}
\begin{itemize}[leftmargin=*, label={\color{columbiablue_official}\tiny$\square$}, parsep=2pt, itemsep=2pt, topsep=0pt]
    \item \textbf{1.1 Context:} Rise of zoning opposition and NIMBYism; protest risks for urban planning.
    \item \textbf{1.2 Case Study:} Austin's housing market, rapid growth, and history of zoning protests.
    \item \textbf{1.3 Research Objectives:} Forecasting per-parcel protest risks via DDPM and LogReg.
\end{itemize}

\section*{2. Literature Review (5-7 pages)}
\begin{itemize}[leftmargin=*, label={\color{columbiablue_official}\tiny$\square$}, parsep=2pt, itemsep=2pt, topsep=0pt]
    \item \textbf{2.1 Zoning Opposition Dynamics:} NIMBYism, homevoter theory, neighborhood resistance.
    \item \textbf{2.2 Quantitative Urban Planning:} Predictive modeling applied to urban phenomena.
    \item \textbf{2.3 Generative Causal Models:} DDPM and Invariant Causal Prediction in tabular models.
\end{itemize}

\section*{3. Mixed-Methods Data and Methodology (8-10 pages) \textcolor{columbiablue_official}{\checkmark}}
\begin{itemize}[leftmargin=*, label={\color{columbiablue_official}\tiny$\boxtimes$}, parsep=2pt, itemsep=2pt, topsep=0pt]
    \item \textbf{3.1 Quantitative Sources \& Panel:} Building the balanced property-year panel (282k parcels).
    \item \textbf{3.2 Causal Prediction \& Leakage:} Preventing temporal leakage across distinct training years.
    \item \textbf{3.3 Qualitative Stakeholder Interviews:} Methodology and thematic coding of 10 interviews.
\end{itemize}

\section*{4. Predictive Modeling Approach (Quantitative) (5-8 pages) \textcolor{columbiablue_official}{\checkmark}}
\begin{itemize}[leftmargin=*, label={\color{columbiablue_official}\tiny$\boxtimes$}, parsep=2pt, itemsep=2pt, topsep=0pt]
    \item \textbf{4.1 Baseline Models:} Setting up Logistic Regression and XGBoost benchmarks.
    \item \textbf{4.2 Conditional Diffusion (DDPM):} Training the generative model for tabular protest risk.
    \item \textbf{4.3 Evaluation Metrics:} RMSE, MAE, CRPS, Scenario Std, out-of-sample validation.
\end{itemize}

\section*{5. Quantitative Results and Counterfactuals (8-10 pages) \textcolor{columbiablue_official}{\checkmark}}
\begin{itemize}[leftmargin=*, label={\color{columbiablue_official}\tiny$\boxtimes$}, parsep=2pt, itemsep=2pt, topsep=0pt]
    \item \textbf{5.1 Baseline Performance:} Feature importance and predictive accuracy of traditional models.
    \item \textbf{5.2 Generative Models:} Evaluating DDPM accuracy and its realistic counterfactuals.
    \item \textbf{5.3 Comparative Analysis \& Policy:} Using evaluated models to inform proactive rezoning.
\end{itemize}

\section*{6. Qualitative Results: Stakeholder Perspectives (8-10 pages)}
\begin{itemize}[leftmargin=*, label={\color{columbiablue_official}\tiny$\square$}, parsep=2pt, itemsep=2pt, topsep=0pt]
    \item \textbf{6.1 Planners and Advocates:} How institutional actors view algorithmic planning tools.
    \item \textbf{6.2 Neighborhood Leaders:} The tension between predictive efficiency and civic participation.
    \item \textbf{6.3 Democratic Implications:} Synthesizing the boundaries of acceptable use.
\end{itemize}

\section*{7. Synthesis and Conclusion (4-6 pages)}
\begin{itemize}[leftmargin=*, label={\color{columbiablue_official}\tiny$\square$}, parsep=2pt, itemsep=2pt, topsep=0pt]
    \item \textbf{7.1 Integrating Findings:} Reconciling models with stakeholder demands for transparency.
    \item \textbf{7.2 Limitations:} ID crosswalk mismatch rates, assumption constraints, model bounding.
    \item \textbf{7.3 Future Work:} Scaling to other jurisdictions; NLP-parsed qualitative petition data.
\end{itemize}

\vspace{1em}
\noindent\textit{Total Estimated Length:} 41-56 pages (excluding Abstract, References, Appendices)

\end{document}
